% ═══════════════════════════════════════════════════════════════════════
\section{Isométries et similitudes}
% ═══════════════════════════════════════════════════════════════════════

\begin{frame}{On descend d'un cran dans la hiérarchie}
  \small Parmi les transformations affines, celles qui intéressent le
  secondaire sont celles qui respectent la \emph{forme} des figures.

  \begin{Enonce}{Définition 1.11 --- Similitude, isométrie, contraction}
    $w$ affine est une \defi{similitude} de rapport $k>0$ si
    \[
      \norme{w(P)-w(Q)} = k\,\norme{P-Q} \qquad \forall P,Q.
    \]
    $k=1$ : \defi{isométrie}. $0<k<1$ : \defi{contraction}.
  \end{Enonce}
  \small C'est la définition de « figures semblables » du secondaire,
  énoncée comme propriété de la \emph{transformation}, non d'un couple
  de figures.
\end{frame}

\begin{frame}{Propositions 1.8 et 1.9}
  \begin{Enonce}{Proposition 1.8 --- Il suffit de regarder la partie linéaire}
    \small $w$ affine de partie linéaire $A$ est une similitude de
    rapport $k$ ssi $\col{x}{y}\mapsto A\col{x}{y}$ en est une, même
    rapport. (La translation disparaît dans la différence
    $w(P)-w(Q)$.)
  \end{Enonce}
  \begin{Enonce}{Proposition 1.9 --- Une composition de similitudes est une similitude}
    \small $w_1,w_2$ similitudes de rapports $k_1,k_2$
    $\Rightarrow$ $w_2\circ w_1$ similitude de rapport $k_1k_2$.
  \end{Enonce}
  \vspace{4pt}
  \textbf{Conséquence.} Les similitudes forment un sous-groupe $\Sim$
  de $\Aff$, les isométries un sous-groupe $\Iso$ de $\Sim$ : la figure
  1.3 est maintenant justifiée.
\end{frame}

\subsection{L'anatomie des similitudes}
\begin{frame}{L'anatomie des similitudes}
  \small Homothéties, rotations, réflexions, translations sont des
  similitudes, donc toutes leurs compositions aussi. On \textbf{admet}
  la réciproque : ce sont les seules --- la partie linéaire d'une
  similitude de rapport $k$ est toujours $k\,r_\theta$ (directe) ou
  $k\,R_\theta$ (indirecte), pour un certain $\theta$. (Démonstration
  complète hors cadre du cours.)
\end{frame}

\begin{frame}{Propositions 1.10 et 1.11}
  \begin{Enonce}{Proposition 1.10 --- Déterminant d'une similitude}
    \small $w$ similitude de rapport $k$, partie linéaire $A$ :
    $\dete(A) = k^2$ ou $-k^2$. Donc $\lvert\dete(A)\rvert = k^2$, et
    $w$ est une contraction ssi $\lvert\dete(A)\rvert < 1$.
  \end{Enonce}
  \begin{Enonce}{Proposition 1.11 --- Aires et rapport de similitude}
    Si $w$ est une similitude de rapport $k$ et $D$ une région du
    plan : $\Aire(w(D)) = k^2\,\Aire(D)$.
  \end{Enonce}
  \begin{Classe}{Le fameux $k^2$}
    \small « Si on double les dimensions, l'aire $\times 4$ » --- l'un
    des résultats les plus mal compris du secondaire. Deux preuves : le
    déterminant (prop.~1.4), ou le fait qu'une aire est un produit de
    deux longueurs, chacune $\times k$. Dimension 3 : $k^3$ pour les
    volumes.
  \end{Classe}
\end{frame}

\begin{frame}{Exemple 1.9 --- Reconnaître une similitude}
  \small $A=\mat{0{,}8}{0}{a}{0{,}6}$ n'est \textbf{jamais} une
  similitude ($\norme{A\vec e_1}\geq 0{,}8 \neq 0{,}6=\norme{A\vec
  e_2}$, contradiction).

  \vspace{6pt}
  $A=\mat{0{,}5}{\alpha}{\beta}{0{,}5}$ : similitude pour
  $\alpha=\beta=0$ (homothétie $\tfrac12$) et pour
  $\alpha=\tfrac{\sqrt3}2,\beta=-\tfrac{\sqrt3}2$ --- on reconnaît la
  matrice $r_{-\pi/3}$ : \textbf{c'est même une isométrie} ($k=1$).
\end{frame}

\subsection{Classification des isométries du plan}
\begin{frame}{Théorème 1.12 --- Les quatre isométries du plan}
  \begin{Enonce}{Théorème 1.12 (admis sans démonstration)}
    \small Toute isométrie du plan est :
    \begin{enumerate}
      \item une \defi{translation} (identité incluse) ;
      \item une \defi{rotation} autour d'un point ;
      \item une \defi{réflexion} par rapport à une droite ;
      \item une \defi{réflexion glissée} : réflexion suivie d'une
        translation parallèle à l'axe.
    \end{enumerate}
  \end{Enonce}
\end{frame}

\begin{frame}{Table 1.2 --- Les isométries, résumé}
  \begin{center}
  \begin{tabular}{llll}
    \toprule
    Isométrie & $\dete(A)$ & Points fixes & Nom au secondaire\\
    \midrule
    Identité         & $+1$ & tout le plan  & --- \\
    Translation      & $+1$ & aucun         & translation \\
    Rotation         & $+1$ & un seul point & rotation \\
    Réflexion        & $-1$ & une droite    & réflexion \\
    Réflexion glissée & $-1$ & aucun         & (absente du programme) \\
    \bottomrule
  \end{tabular}
  \end{center}
\end{frame}

\begin{frame}{Pour la classe 1.7 --- La grande absente}
  \begin{Classe}{La grande absente}
    \small La réflexion glissée ne figure pas au programme québécois.
    Elle explique les traces de pas dans la neige, la disposition des
    feuilles, plusieurs frises. La liste « translation, rotation,
    réflexion » que retiennent les élèves est \textbf{incomplète} :
    composer une réflexion et une translation donne, en général,
    quelque chose qui n'est aucune des trois.
  \end{Classe}
  \FigBox[0.3\textheight]{%
    \begin{tikzpicture}[scale=0.85]
      \draw[droite] (-0.5,0) -- (8.2,0);
      \foreach \i in {0,1,2,3}{
        \begin{scope}[shift={(2*\i,0)}, scale=0.55]
          \fill[fig] \motifF;
        \end{scope}
        \begin{scope}[shift={(2*\i+1,0)}, scale=0.55, yscale=-1]
          \fill[figim] \motifF;
        \end{scope}}
      \node[etiq, right] at (8.3,0) {axe};
    \end{tikzpicture}}
  \begin{center}\small Figure 1.9 --- Une frise engendrée par une
  réflexion glissée.\end{center}
\end{frame}

\subsection{Tout se ramène aux réflexions}
\begin{frame}{Théorème 1.13 --- Composée de deux réflexions}
  \begin{Enonce}{Théorème 1.13}
    \small Soient $d_1,d_2$, $R_1,R_2$ les réflexions correspondantes.
    \begin{enumerate}
      \item Si $d_1\parallel d_2$ : $R_2\circ R_1$ = translation de
        vecteur $2\vec v$, $\vec v$ perpendiculaire de $d_1$ à $d_2$.
      \item Si $d_1,d_2$ se coupent en $M$, angle $\varphi$ :
        $R_2\circ R_1$ = rotation de centre $M$, angle $2\varphi$.
    \end{enumerate}
  \end{Enonce}
  \FigBox[0.4\textheight]{%
    \begin{tikzpicture}[scale=1.25]
      \draw[droite] (-1.8,0) -- (1.8,0) node[etiq, right, black] {$d_1$};
      \draw[droite] (35:-1.8) -- (35:1.8) node[etiq, right, black] {$d_2$};
      \node[pt, label={[etiq]below left:$M$}] at (0,0) {};
      \draw[sectioncolor, thin] (0.6,0) arc (0:35:0.6);
      \node[etiq, sectioncolor] at (17:0.82) {$\varphi$};
      \begin{scope}[shift={(0.75,0.25)}, scale=0.5]
        \fill[fig] \motifF;
      \end{scope}
      \begin{scope}[cm={1,0,0,-1,(0,0)}]
        \begin{scope}[shift={(0.75,0.25)}, scale=0.5]
          \fill[figaux] \motifF;
        \end{scope}
      \end{scope}
      \begin{scope}[rotate=70]
        \begin{scope}[shift={(0.75,0.25)}, scale=0.5]
          \fill[figim] \motifF;
        \end{scope}
      \end{scope}
      \draw[-{Stealth[length=2mm]}, firebrick3, thick]
        (1.35,0.75) arc (25:95:1.5);
      \node[etiq, firebrick3] at (60:1.85) {$2\varphi$};
    \end{tikzpicture}}
  \begin{center}\small Figure 1.10\end{center}
\end{frame}

\begin{frame}{Corollaire 1.14, et une leçon de didactique}
  \begin{AuTableau}[Idée de la démonstration --- au tableau]
    \small Cas $2$ : quitte à translater, $M=\vec 0$ ; par
    $R_{\alpha+\varphi}R_\alpha$, développement trigonométrique
    $\Rightarrow r_{2\varphi}$.
  \end{AuTableau}
  \begin{Enonce}{Corollaire 1.14 --- Trois réflexions suffisent}
    Toute isométrie du plan est la composée d'au plus \textbf{trois}
    réflexions.
  \end{Enonce}
  \begin{Classe}{Une leçon de didactique}
    \small La réflexion n'est pas une transformation « parmi
    d'autres » : c'est \emph{la} transformation fondamentale, qui
    engendre toutes les autres. On peut faire vivre ce résultat avec
    deux miroirs, ou deux pliages successifs : les élèves découvrent
    expérimentalement que deux réflexions donnent une rotation d'angle
    double de l'angle entre les plis.
  \end{Classe}
\end{frame}
